\documentclass{article}
\usepackage[UTF8]{ctex}
\usepackage{amsmath,amssymb,amsfonts}
\usepackage{enumitem}
\usepackage{geometry}
\geometry{a4paper, margin=1in}
% === 错题本排版优化（不参与题目标识） ===
\setlength{\parindent}{0pt}
\setlength{\parskip}{0.35em plus 0.10em minus 0.08em}
\setlength{\emergencystretch}{3em}
\setlist{itemsep=0.25em, topsep=0.40em, parsep=0pt}
\allowdisplaybreaks[2]
% === 错题本排版优化结束 ===
\begin{document}
\section*{中国石油大学(北京)2025年高等代数真题}
\section*{27．中国石油大学（北京）2025年高等代数试卷}
一．填空题．每题4分，共24分．\\

\subsection{题目1}
1．已知 \(A=\left(a_{i j}\right)\) 是实数域上三阶非零矩阵，\(A_{i j}\) 是元素 \(a_{i j}\) 在行列式 \(|A|\) 中的代数余子式，若\(A_{i j}+a_{i j}=0(i, j=1,2,3)\) ，则 \(|A|=\) \(\_\_\_\_\)．

\subsection{题目2}
2．设 \(A, B\) 均为三阶实矩阵，\(D^{\mathrm{T}}\) 为 \(D\) 的转置。

\[
D^{\mathrm{T}} A D=\left(\begin{array}{lll}
3 & 2 & 0 \\
2 & 1 & 0 \\
0 & 0 & 2
\end{array}\right)
\]

若 \(D=\left(y_{1}, y_{2}, y_{3}\right), C=\left(y_{1}, 2 y_{2}, y_{3}-y_{1}\right)\) ，则 \(C^{\mathrm{T}} A C=\) \(\_\_\_\_\)．

\subsection{题目3}
3．已知 \(V=\left\{\left.\left(\begin{array}{cc}b+b \mathrm{i} & c \\ a & b-b \mathrm{i}\end{array}\right) \right\rvert\, a, b, c \in \mathbb{R}\right\}\) 是实数域上的线性空间，则 \(V\) 的维数为 \(\_\_\_\_\)．\\
\(\_\_\_\_\)

\subsection{题目4}
4．已知实矩阵 \(A\) 的特征多项式为 \(f(\lambda)=(\lambda+1)^{2}(\lambda-5)^{3}\) ，其最小多项式为 \(g(\lambda)=(\lambda+1)^{2}(\lambda-5)^{2}\)则 \(A\) 的 Jordan 标准形为 \(\_\_\_\_\)．

\subsection{题目5}
5．三阶实对称矩阵 \(A\) 的特征值为 \(\lambda_{1}=1, \lambda_{2}=-1, \lambda_{3}=0, \lambda_{1}\) 的特征向量为 \(p_{1}=(1, x, 1)^{\mathrm{T}}, \lambda_{3}\) 的特征向量为 \(p_{3}=(x, x+1,1)^{\mathrm{T}}, x \in \mathbb{R}\) ，则 \(x=\) \(\_\_\_\_\)，\(\lambda_{2}\) 的特征向量为 \(p_{2}=\) \(\_\_\_\_\)．

\subsection{题目6}
6．\(e_{1}, e_{2}, e_{3}, e_{4}, e_{5}\) 为欧氏空间 \(V\) 上的一组标准正交基，\(W\) 是由向量组 \(e_{3}, e_{1}+e_{3}, e_{1}+e_{2}+e_{3}\) 生成的子空间，则它的正交补 \(W^{\perp}=\) \(\_\_\_\_\)。

二、选择题．1－4 为单选，5－6 为多选，多选题多选，少选，错选均不得分，每题6分，共36分。\\

\subsection{题目7}
1．已知实对称矩阵

\[
A=\left(\begin{array}{cccc}
a_{11} & a_{12} & a_{13} & a_{14} \\
a_{21} & a_{22} & a_{23} & a_{24} \\
a_{31} & a_{32} & a_{33} & a_{34} \\
a_{41} & a_{42} & a_{43} & a_{44}
\end{array}\right)
\]

满足 \(A^{8}+9 A^{5}+3 A^{4}+A^{2}+4 E=O\) ，则 \(\left|\begin{array}{ccc}a_{11} & a_{12} & a_{14} \\ a_{12} & a_{22} & a_{24} \\ a_{14} & a_{24} & a_{44}\end{array}\right|\) 必然（）\\
A．＞0；B．＝0；C．＜0；D．不确定。

\subsection{题目8}
2．下列多项式在复数域上有重根的是（）\\
A．\(x^{n}+1\) ；\\
B． \(1+x+\frac{x^{2}}{2!}+\cdots+\frac{x^{n}}{n!}\) ；\\
C．\(n x^{n+1}-(n+1) x^{n}+1\) ；\\
D．\(x^{n}+x^{n-1}+\cdots+1\) ．

\subsection{题目9}
3．实矩阵 \(A=\left(\begin{array}{ccc}2 & 4 & -2 \\ 4 & 2 & 2 \\ -2 & 2 & 5\end{array}\right), B=\left(\begin{array}{ccc}1 & 0 & 0 \\ 0 & -2 & 0 \\ 0 & 0 & 4\end{array}\right)\) ，则 \(A\) 与 \(B()\)\\
A．合同，相似；B．不合同，相似；C．合同，不相似；D．不合同，不相似．

\subsection{题目10}
4．\(A=\left(a_{1}, a_{2}, a_{3}, a_{4}\right)\) 是 \(A\) 阶正交阵，\(B=\left(\begin{array}{l}\alpha_{1}^{\mathrm{T}} \\ \alpha_{2}^{\mathrm{T}} \\ \alpha_{3}^{\mathrm{T}}\end{array}\right), \beta=\left(\begin{array}{l}1 \\ 1 \\ 1\end{array}\right), k, l\) 为任意实数，则线性方程组\(B X=\beta\) 通解为（）\\
A．\(k a_{1}+l a_{2}+a_{2}+a_{3}+a_{4}\) ．\\
B．\(k a_{4}+a_{1}+a_{2}+a_{3}\) ．\\
C．\(k a_{3}+l a_{4}+a_{1}+a_{2}+a_{3}\) ．\\
D．\(k a_{2}+a_{1}+a_{3}+a_{4}\) ．

\subsection{题目11}
5．不是实数域上线性空间 \(\mathbb{R}^{n}\) 的子空为（）

\begin{itemize}
  \item[A．] \(V_{1}=\left\{\left(x_{1}, x_{2}, \cdots, x_{n}\right)^{\mathrm{T}} \mid x_{1}+x_{3}+x_{n}=0, \forall x_{i} \in \mathbb{R}\right\}\) ．
  \item[B．] \(V_{2}=\left\{X=k_{1} \xi_{1}+2 \xi_{2} \mid \forall k \in \mathbb{R}\right\}\) ，其中 \(\xi_{1}, \xi_{2} \in \mathbb{R}^{n}\) 线性无关。
  \item[C．] \(V_{3}=\left\{\left(x_{1}, x_{2}, \cdots, x_{n}\right)^{\mathrm{T}} \mid x_{1}^{2}+x_{2}^{2}+\cdots+x_{n}^{2}=1, \forall x_{i} \in \mathbb{R}\right\}\) ．
  \item[D．] \(V_{4}=\left\{\left(x_{1}, x_{2}, \cdots, x_{n}\right)^{\mathrm{T}} \mid x_{1} x_{2} \cdots x_{n}=0, \forall x_{i} \in \mathbb{R}\right\}\) ．
  \item[E．] \(V_{5}=\left\{\left(x_{1}, x_{2}, \cdots, x_{n}\right)^{\mathrm{T}} \mid x_{1}^{2}+x_{2}^{2}+\cdots+x_{n}^{2}=0, \forall x_{i} \in \mathbb{R}\right\}\) ．
\end{itemize}

\subsection{题目12}
6．下列正确的是（）

\begin{itemize}
  \item[A．] \(\alpha, \beta, y\) 为欧氏空间 \(V\) 中的向量，且 \(\alpha, \beta, y\) 正交，则对任意的 \(k, l, m \in \mathbb{R}, k \alpha\) 与 \(l \beta+m y\) 正交。
  \item[B．] \(n\) 维欧氏空间 \(V\) 中存在 \(n+1\) 个两两正交的非零向量。
  \item[C．] \(A\) 为正交阵，则 \(-\left(A^{\mathrm{T}}\right)^{-1}\) 也为正交阵。
  \item[D．] 欧氏空间上保持向量距离不变的线性变换一定是正交变换．
  \item[E．] 欧氏空间上保持向量夹角不变的线性变换一定是正交变换．
\end{itemize}
\end{document}
